\documentclass[10pt,a4paper]{amsart}
\usepackage[margin=19mm]{geometry}
\usepackage{amsmath,amssymb,mathtools}
\usepackage[scr=rsfs,cal=euler]{mathalfa}
\usepackage{xfrac}
\usepackage[unicode,hidelinks,bookmarks=false]{hyperref}
\newcommand{\ie}{i.e.\ }
\newcommand{\cf}{cf.\ }
\newcommand{\resp}{respectively\ }
\newcommand{\Subsection}{\subsection}
\providecommand{\nobreakdash}{\nobreak\hbox{-}}
\setlength{\emergencystretch}{2em}
\multlinegap0pt
\usepackage{bm}
\usepackage{bbm}
\usepackage{dsfont}
\usepackage{xspace}
\usepackage{cancel}
\usepackage{mathtools}
\usepackage{amsthm}
\makeatletter
\@ifundefined{thm}{\newtheorem{thm}{Thm}}{}
\@ifundefined{lemma}{\newtheorem{lemma}[thm]{Lemma}}{}
\makeatother
\newcommand{\II}{\mathrm{II}}
\newcommand{\I}{\mathrm{I}}
\newcommand{\N}{\mathbb{N}}
\newcommand{\abs}[1]{\left|#1\right|}
\newcommand{\brkt}[1]{\left(#1\right)}
\newcommand{\dd}{\mathrm{d}}
\title[Multipliers for Hardy-Orlicz spaces and applications]{Multipliers for Hardy-Orlicz spaces and applications}
\author{Odysseas Bakas, Sandra Pott, Salvador Rodriguez-Lopez, Alan Sola}
\date{Source: arXiv:2306.09874v1 (CC BY 4.0)}
\begin{document}
\maketitle

\noindent\emph{Source and licence.} Multipliers for Hardy-Orlicz spaces and applications. Version arXiv:2306.09874v1 (CC BY 4.0), distributed under the Creative Commons Attribution 4.0 International licence, as recorded in the arXiv licence metadata for this exact version and shown on the abstract page https://arxiv.org/abs/2306.09874v1. Full rights notice: licenses/arxiv-2306.09874-CC-BY-4.0.txt.\par\smallskip

\noindent\textit{Editorial scope.} Counted unit: the Lemma whose source label is \texttt{lem:uncond1} in the article (source0.tex lines 865--881, source label \texttt{lem:uncond1}), delivered together with the article's own complete original proof of it (source0.tex lines 882--996, one \texttt{proof} environment of 115 lines). No mathematics is written, repaired or re-derived: statement and proof are the article's own text line for line. Required context retained uncounted: none. Every definition, display and label of those lines is retained unchanged. Omitted from this package and not redistributed: 1--864, 997--2388. Nothing inside a retained excerpt is omitted or altered: every retained body is delivered exactly as the article's lines are written, so no omission receipt of any kind accompanies this package. Citations: the retained excerpts carry no citation command at all, so no bibliography entry of the article is transcribed and no reference is redistributed. Reference adaptation: none was needed and none was made; every reference inside the delivered unit resolves inside it, so no unresolved reference or citation is produced. Printed slips kept literal: no printed word, symbol, hypothesis or number of the counted statement or of its proof was altered. Layout and class: the article's own class is \texttt{amsart} with packages including amssymb, amsmath, amsthm, amsrefs, hyperref, mathrsfs, mathabx, enumitem, esint, todonotes, xcolor; the wrapper is the lane's pinned \texttt{amsart} editorial wrapper at 10pt on A4, which restates the theorem-like environments the retained text needs and adds the article's own macro definitions that the retained text needs (\texttt{\textbackslash I}, \texttt{\textbackslash II}, \texttt{\textbackslash N}, \texttt{\textbackslash abs}, \texttt{\textbackslash brkt}, \texttt{\textbackslash dd}), with extra wrapper packages (bm, bbm, dsfont, xspace, cancel, mathtools). The delivered numbering is the unit's own. Assets: no retained span contains a figure environment, a graphics-inclusion command, a PDF-page inclusion, a table or a TikZ picture; no image file is copied into this package, the package declares no asset dependency and no image file is redistributed with it. Licence: version arXiv:2306.09874v1 of this article is distributed under the Creative Commons Attribution 4.0 International licence, as recorded in the arXiv licence metadata for this exact version and shown on the abstract page https://arxiv.org/abs/2306.09874v1. A full rights notice is delivered at licenses/arxiv-2306.09874-CC-BY-4.0.txt. Characters: every retained excerpt is pure ASCII, so the delivered engine, XeLaTeX with its default Latin Modern text font, reports no missing character.
\begin{lemma}\label{lem:uncond1}
If $n \in \N_0$, %
then for every interval $I$ with $\abs{I}= 2^{-n}$ and for every analytic trigonometric polynomial $f$, 
\begin{enumerate}
%
\item \label{geometric}
\[    \int_I \left| f( e^{i 2 \pi \theta} ) \right|^2 \, \dd \theta \lesssim  \abs{I}  \brkt{\sum_{j=0}^{2^{n}-1} \abs{f_j}}^2 +
    \abs{I}  \sum_{k=0}^\infty \brkt{\sum_{j=2^k 2^n}^{2^{k+1} 2^{n}-1} \abs{f_j}}^2,
\]
\item \label{arithmetic}
\[    \int_I \left| f( e^{i 2 \pi \theta} ) \right|^2 \, \dd \theta \lesssim 
    \abs{I}  \sum_{k=0}^\infty \brkt{\sum_{j=k2^n}^{(k+1)2^{n}-1} \abs{f_j}}^2,
\]
\end{enumerate}
where $f_j := \widehat{f} (j)$, $j \in \N_0$. 
%
\end{lemma}

\begin{proof} Let $|I| = 2^{-n}$.
We remark that Part (\ref{geometric}) of the lemma follows from Part (\ref{arithmetic}) by a simple $\ell^1 - \ell^2$ estimate, so it suffices to prove Part (\ref{arithmetic}).

Note that
\begin{align*}
\int_I f(e^{i 2 \pi \theta})\overline{f(e^{i 2 \pi \theta})} \, \dd \theta & = \sum_{k_1 = 0}^\infty\sum_{k_2=0}^\infty \sum_{j=k_1 2^{n}}^{(k_1+1)2^{n}-1}\sum_{l=k_2 2^n}^{(k_2 +1) 2^{n}-1} f_j \overline{f_l} \int_I e ^{i 2\pi (j-l) \theta} \, \dd \theta \\
&=  \I + \II,
\end{align*}
where
\[
\I := \sum_{\substack{ k_1, k_2 \geq 0:\\ |k_1-k_2| \leq 1}} \sum_{j=k_1 2^{n}}^{(k_1+1)2^{n}-1}\sum_{l=k_2 2^n}^{(k_2 +1) 2^{n}-1} f_j \overline{f_l} \int_I e^{i 2\pi (j-l) \theta} \, \dd \theta
\]
and
\[
\II := \sum_{\substack{ k_1, k_2 \geq 0:\\ |k_1-k_2| \geq 2}} \sum_{j=k_1 2^{n}}^{(k_1+1)2^{n}-1}\sum_{l=k_2 2^n}^{(k_2 +1) 2^{n}-1} f_j \overline{f_l} \int_I e^{i 2\pi (j-l) \theta} \, \dd \theta .
\]

Term $\I$ may easily be estimated by
\[
| \I | \leq
\sum_{\substack{k_1, k_2 \geq 0:\\ |k_1-k_2| \leq 1}} \sum_{j=k_1 2^{n}}^{(k_1+1)2^{n}-1}\sum_{l=k_2 2^n}^{(k_2 +1) 2^{n}-1} | f_j | \cdot |  f_l | \cdot 2^{-n}
\leq 3 \cdot 2^{-n} \sum_{k \geq 0} \left(\sum_{j=k 2^{n}}^{(k+1)2^{n}-1} |f_j| \right)^2.
\]
For term $\II$, consider first the case $I = [0,2^{-n}]$. Then
\begin{align*}
&\sum_{\substack{k_1, k_2 \geq 0:\\ |k_1-k_2| \geq 2}} \sum_{j=k_1 2^{n}}^{(k_1+1)2^{n}-1}\sum_{l=k_2 2^n}^{(k_2 +1) 2^{n}-1} f_j \overline{f_l} \int_I e^{i 2\pi (j-l) \theta} \, \dd \theta = \\
& \sum_{\substack{k_1, k_2 \geq 0:\\ |k_1-k_2| \geq 2}}
\sum_{j=k_1 2^{n}}^{(k_1+1)2^{n}-1} \sum_{l=k_2 2^n}^{(k_2 +1) 2^{n}-1} f_j \overline{f_l} \, \frac{1 - e^{i 2 \pi (j-l) 2^{-n}}}{i 2\pi (j-l)} =  \\
&  \text{III} + \text{IV},
\end{align*}
where
\[
\text{III} :=  \sum_{\substack{ k_1, k_2 \geq  0:\\ k_1 \neq k_2} }^\infty 
\sum_{j = k_1 2^{n}}^{(k_1+1) 2^{n}-1} 
\sum_{l = k_2 2^{n}}^{(k_2+1) 2^{n}-1} f_{j} \overline{f_{l}} \frac{1}{i 2\pi ( j-l)}
\]
and
\[
\text{IV} := - \sum_{ \substack{k_1, k_2 \geq  0: \\ k_1 \neq k_2 }}^\infty 
\sum_{j = k_1 2^{n}}^{(k_1+1) 2^{n}-1} 
\sum_{l = k_2 2^{n}}^{(k_2+1) 2^{n}-1}  f_{j} \overline{f_{l}}  \frac{  e^{i 2 \pi (j-l) 2^{-n}}}{i 2\pi  ( j-l)} . 
\]

Moreover,
\begin{align*}
\text{III} &= \sum_{\substack{k_1, k_2\geq 0:\\ |k_1-k_2| \geq 2}}^\infty
\sum_{j_1=0}^{2^{n}-1} \sum_{j_2=0}^{2^{n}-1} f_{k_1 2^n +j_1} \overline{f_{k_2 2^n + j_2}}  \frac{1}{i 2\pi ( 2^n (k_1-k_2) + (j_1-j_2))} \\
& = \text{V} + \text{VI},
\end{align*}
where
\[
\text{V}:= \sum_{\substack{k_1, k_2\geq 0:\\ |k_1 -k_2|\geq 2}}
\sum_{j_1=0}^{2^{n}-1} \sum_{j_2=0}^{2^{n}-1} f_{k_1 2^n +j_1} \overline{f_{k_2 2^n + j_2}}  \frac{1}{i 2\pi 2^n (k_1-k_2)} 
\]
and
\begin{multline*}
\text{VI} := \\
\sum_{\substack{ k_1, k_2\geq 0:\\ |k_1-k_2| \geq 2}}^\infty
\sum_{j_1=0}^{2^{n}-1} \sum_{j_2=0}^{2^{n}-1} f_{k_1 2^n +j_1} \overline{f_{k_2 2^n + j_2}}  \frac{j_2 -j_1}{i 2\pi ( 2^n (k_1-k_2) + (j_1-j_2)) 2^n (k_1-k_2) } .
\end{multline*}

Term V can be estimated by
\begin{align*}
& \left|\sum_{|k_1 -k_2|\geq 2}
\sum_{j_1=0}^{2^{n}-1} \sum_{j_2=0}^{2^{n}-1} f_{k_1 2^n +j_1} \overline{f_{k_2 2^n + j_2}}  \frac{1}{i 2\pi 2^n (k_1-k_2)}
\right| = \\
& \frac{2^{-n}}{2 \pi} \left|\sum_{|k_1 -k_2|\geq 2}
\left(\sum_{j_1=0}^{2^{n}-1}  f_{k_1 2^n +j_1} \right)
\left( \sum_{j_2=0}^{2^{n}-1} \overline{f_{k_2 2^n + j_2}} \right) \frac{1}{k_1-k_2}
\right|
\lesssim \\
&  2^{-n} \sum_{k \geq 1} \left(\sum_{j=k 2^{n}}^{(k+1)2^{n}-1} |f_j| \right)^2,
\end{align*}
where we use the boundedness of the discrete Hilbert transform on $l^2$.

For term $\text{VI}$, we observe that
\[
\left|\frac{j_2 -j_1}{i 2\pi ( 2^n (k_1-k_2) + (j_1-j_2))  2^n (k_1-k_2)}
 \right| \leq 2 \cdot 2^{-n} \frac{1}{2\pi  (k_1-k_2)^2}
\]
for $|k_1-k_2| \geq 2$ and $j_1, j_2 \in \{ 0, 1, \cdots, 2^n-1\}$, which yields
\begin{align*}
 & \left|\sum_{\substack{ k_1, k_2\geq 0:\\ |k_1-k_2| \geq 2}}^\infty
\sum_{j_1=0}^{2^{n}-1} \sum_{j_2=0}^{2^{n}-1} f_{k_1 2^n +j_1} \overline{f_{k_2 2^n + j_2}}  \frac{j_2 -j_1}{i 2\pi ( 2^n (k_1-k_2) + (j_1-j_2)) 2^n (k_1-k_2)   } \right| \leq \\
&   \frac{2^{-n}}{ \pi} \sum_{|k_1 -k_2|\geq 2}
\left(\sum_{j_1=0}^{2^{n}-1}  |f_{k_1 2^n +j_1} |\right)
\left( \sum_{j_2=0}^{2^{n}-1} |f_{k_2 2^n + j_2}| \right) \frac{1}{(k_1-k_2)^2}
\lesssim \\
& 2^{-n} \sum_{k \geq 0} \left(\sum_{j=k 2^{n}}^{(k+1)2^{n}-1} |f_j| \right)^2,
\end{align*}
where we use the the boundedness of the linear operator on $\ell^2$ induced by the matrix
\[
\left(\frac{1}{(k_1-k_2)^2}\right)_{k_1, k_2 \geq 0, |k_1 -k_2|\geq 2},
\]
as can be verified e.g. by appealing to the Schur test.

Note that we have yet to estimate term IV. Instead of doing this directly, we observe that for $I = [a, a+ 2^{-n}]$, term II becomes
\[ 
\sum_{\substack{k_1, k_2 \geq 0:\\ |k_1-k_2| \geq 2}}
\sum_{j=k_1 2^{n}}^{(k_1+1)2^{n}-1} \sum_{l=k_2 2^n}^{(k_2 +1) 2^{n}-1} f_j \overline{f_l} \, 
\frac{e^{i 2 \pi (j-l)a} - e^{i 2 \pi (j-l) (a+2^{-n})}}{i 2\pi (j-l)} = \text{VII} + \text{VIII},
\]
where
\[
\text{VII} := \sum_{\substack{k_1, k_2 \geq 0:\\ |k_1-k_2| \geq 2}}
\sum_{j=k_1 2^{n}}^{(k_1+1)2^{n}-1} \sum_{l=k_2 2^n}^{(k_2 +1) 2^{n}-1} e^{i 2 \pi j a} f_j \overline{e^{i 2 \pi l a} f_l} \, 
    \frac{1}{i 2\pi (j-l)}
\]
and
\[
\text{VIII} := - \sum_{\substack{k_1, k_2 \geq 0:\\ |k_1-k_2| \geq 2}}
\sum_{j=k_1 2^{n}}^{(k_1+1)2^{n}-1} \sum_{l=k_2 2^n}^{(k_2 +1) 2^{n}-1} e^{i 2 \pi j (a + 2^{-n})} f_j \overline{e^{i 2 \pi l (a+ 2^{-n})} f_l} \, 
\frac{1}{i 2\pi (j-l)} .
\]
To handle these last expressions, note that $\text{VII}$ and $\text{VIII}$ have exactly the same form as term $\text{III}$, just with unimodular factors on the Fourier coefficients $f_j$. This observation combined with our previous estimates finishes the proof of the lemma. \end{proof}

\end{document}
